\input format.tex

\usepackage{graphicx}
\graphicspath{{cores/}}

\begin{document}

%\vspace*{8mm}
\parindent0pt
%% 各章节
\setlength{\arrayrulewidth}{1pt}
\fontsize{9.3pt}{11pt}\selectfont
\color{gray2}

{\bf\sanhao 检测结果总览}

\vspace*{2mm}

\begin{LRaside}{肠道菌群概况}
\noindent\centering
\includegraphics[width=\linewidth]{gaikuang_green.pdf}
\jiuhao
\noindent\begin{tabular}{@{}m{0.68cm}<{\centering}@{}m{0.68cm}<{\centering}@{}m{0.68cm}<{\centering}@{}m{0.68cm}<{\centering}@{}}
\includegraphics[width=.41cm]{sun_green.pdf} & \tikz\draw[gray,fill=gray](0,0) circle[radius=4.1pt]; & \tikz\draw[gray,fill=gray](0,0) circle[radius=4.1pt]; & \tikz\draw[gray,fill=gray](0,0) circle[radius=4.1pt]; \\[-4pt]
\color{gray2}当次 & \color{gray2} & \color{gray2} & \color{gray2}
\\
\end{tabular}
\asidebreak % 分栏
1、您本次肠道菌群共检测出细菌种类169种，肠道菌群多样性
正常，说明您的肠道菌群组成较为丰富，菌群失调风险较低。\par\bigskip
2、您肠道中的有益菌含量
正常，高于93.75{\%}的参考人群。\par\bigskip
3、您肠道内的有害菌含量
与参考人群相当。
\end{LRaside}

\begin{LRaside}{肠道菌群含量}
\noindent
\includegraphics[width=\linewidth]{fenbu_orange.pdf}

\jiuhao
\noindent\begin{tabular}{@{}m{0.68cm}<{\centering}@{}m{0.68cm}<{\centering}@{}m{0.68cm}<{\centering}@{}m{0.68cm}<{\centering}@{}}
\includegraphics[width=.41cm]{sun_orange.pdf} & \tikz\draw[gray,fill=gray](0,0) circle[radius=4.1pt]; & \tikz\draw[gray,fill=gray](0,0) circle[radius=4.1pt]; & \tikz\draw[gray,fill=gray](0,0) circle[radius=4.1pt]; \\[-4pt]
\color{gray2}当次 & \color{gray2} & \color{gray2} & \color{gray2}
\\
\end{tabular}
\asidebreak % 分栏
1、您肠道内的
常见肠道菌含量均处于正常范围，共同维持肠道的正常运转。\par\bigskip
2、您的肠道内
检测出7种致病菌，含量均未超标，但在您免疫力低下时仍可能会损害您的肠道健康。
\end{LRaside}


\begin{LRaside}{肠道菌群营养功能分析}
\noindent
\includegraphics[width=\linewidth]{gongneng_green.pdf}

\jiuhao
\noindent\begin{tabular}{@{}m{0.68cm}<{\centering}@{}m{0.68cm}<{\centering}@{}m{0.68cm}<{\centering}@{}m{0.68cm}<{\centering}@{}}
\includegraphics[width=.41cm]{sun_green.pdf} & \tikz\draw[gray,fill=gray](0,0) circle[radius=4.1pt]; & \tikz\draw[gray,fill=gray](0,0) circle[radius=4.1pt]; & \tikz\draw[gray,fill=gray](0,0) circle[radius=4.1pt]; \\[-4pt]
\color{gray2}当次 & \color{gray2} & \color{gray2} & \color{gray2}
\\
\end{tabular}
\asidebreak %分栏
您肠道内参与合成B族维生素和类固醇等营养物质的菌群含量均处于正常范围，帮助您维持机体正常生理功能。
\end{LRaside}

颜色指示 \includegraphics[scale=1]{greencircle.pdf} {\bf 绿色表示健康}   \includegraphics[scale=1]{orangecircle.pdf}  {\bf 橙色表示需要关注}  \includegraphics[scale=1]{redcircle.pdf} {\bf 红色表示有风险}

\end{document}
